\section{Discussion}
\label{sec:discussion_limitation_future}

With empirical insights gleaned from \method and \dataset, we discuss three key areas for improvement for contemporary LLMs safety and general AI safety research. 

\paragraph{Addressing AI safety comprehensively and openly.}
Due to the scarcity of publicly available safety training resources and lessons, the research community remains to face substantial challenges toward building robustly safe models. The AI community demands publicly shared norms, best practices, and technical standards to identify and quantify unexpected system outputs, and to develop corresponding defenses before risks arise in public settings. By sharing the insights of \method and the release of \dataset, we take concrete steps toward more open conversations around safety training resources and practices. Building on this, we call for similar open engagement for safety resources to address model vulnerabilities comprehensively and concretely.


\paragraph{Upgrading safety evaluation methods, tasks, and metrics to tailor to evolving model capabilities.}
Many existing safety benchmarks are either contaminated \citep{golchin2023time} or saturated \citep{zheng2023lmsyschat1m}, and existing classifiers and metrics can often be inaccurate. Developing robust, safe models requires a dynamic pipeline between exhaustive evaluation (to uncover ever-evolving model vulnerabilities) and safety adaption (to improve identified behaviors). Thus, innovating scalable and evolving safety evaluation approaches is crucial for accurately assessing model safety levels. In particular, ideal safety evaluations should go beyond standard red-teaming approaches that typically involve a small team of experts, only explore a narrow risk domain, and remain static despite facing improved model performances. With \method, we make a substantial attempt at a broader assessment of model shortcomings across wider risk categories and attack types. For the future, we call for upgrades of safety evaluation methods, tasks, and metrics to keep pace with improving model capabilities.

\paragraph{Scrutinizing the training recipes and internal mechanisms that lead to robust safe models.}
This work shows that simple but effective supervised fine-tuning on high-quality safety data can already lead to substantial model safety improvements. However, we need a deep-down understanding of the best practices of safety alignment, \eg pros and cons between SFT vs. DPO vs. PPO; end-to-end safety-trained LMs vs. plug-in safety filters; direct vs. elaborated refusal. In addition, it's valuable to understand when and why a safety alignment approach may or may not work by probing into its underlying mechanisms. Such anatomy of safety alignment may involve disentangling competing learning objectives of safety-trained models, \ie instruction-following vs. refusal \citep{wei2023jailbroken}, and developing novel approaches that better facilitate the Pareto frontier across such abilities (\eg unlearning \citep{goel2024corrective}, contrastive learning \citep{zou2024improving}, representation engineering at neuron or rank level \citep{wei2024assessing}). Finally, it's important to troubleshoot superficial safety alignment processes \citep{zhou2024lima, lubana2023mechanistic} revealed by malicious data poisoning \citep{qi2023fine} or unexpected backdoor behaviors \citep{hubinger2024sleeper}.

